\documentclass[11pt,a4paper]{article}

% ------------------------------------------------------------------ %
% Paquetes (replica el preámbulo del report principal del solver)
% ------------------------------------------------------------------ %
\usepackage[margin=2.2cm]{geometry}
\usepackage[spanish,es-noquoting]{babel}
\usepackage[utf8]{inputenc}
\usepackage{amsmath,amssymb,mathtools}
\usepackage{booktabs}
\usepackage{array}
\usepackage{xcolor}
\usepackage{listings}
\usepackage{siunitx}
\usepackage{microtype}
\usepackage{enumitem}
\usepackage{algorithm}
\usepackage{algpseudocode}
\usepackage[hidelinks]{hyperref}

\lstdefinestyle{cfg}{
  basicstyle=\ttfamily\footnotesize,
  breaklines=true,
  columns=fullflexible,
  frame=single,
  rulecolor=\color{black!30},
  showstringspaces=false,
}
\lstset{style=cfg}

% ------------------------------------------------------------------ %
% Macros de conveniencia
% ------------------------------------------------------------------ %
\newcommand{\Vm}{V_{m}}
\newcommand{\Iion}{I_{\text{ion}}}
\newcommand{\Iext}{I_{\text{ext}}}
\newcommand{\Dten}{\mathbf{D}}
\newcommand{\grad}{\nabla}
\newcommand{\divg}{\nabla\!\cdot}
\newcommand{\code}[1]{\texttt{\small #1}}
\newcommand{\gpode}{\code{gaussPointODE}}
\newcommand{\ccrec}{\code{cellCentredReconstruct}}
\newcommand{\states}{\mathbf{s}}
\newcommand{\Vbar}{\overline{V}_{m}}
\newcommand{\Vbarc}{\overline{V}_{m,c}}
\newcommand{\avg}[1]{\left\langle #1 \right\rangle}
\newcommand{\Omegac}{\Omega_{c}}

\title{Por qué los \emph{states} de \ccrec\ fallan en un frente TNNP propagante\\
       mientras \gpode\ funciona\\[2pt]
       \large Un diagnóstico de la falla de ignición y de conducción en
       \texttt{highOrderElectroActivationFoamImplicitPETSc}}
\author{Nota de diagnóstico del solver}
\date{\today}

\begin{document}
\maketitle

\begin{abstract}
El solver ofrece dos formas de manejar las variables de estado iónicas
(\emph{states}) $\states=(s_1,\dots,s_{N})$ en la cuadratura de alto orden de la
corriente iónica: \gpode\ (el camino heredado, una integración del ODE rígido por
punto de Gauss) y \ccrec\ (un ODE por centro de celda, y los \emph{states} luego
reconstruidos en los puntos de Gauss por una LRE de alto orden de orden
\code{p1}/\code{p2}/\code{p3}). En el benchmark de
Niederer\,\emph{et al.}~(2012) con el modelo celular TNNP, \ccrec\ no logra
producir una activación propagante a medida que se refina el paso de tiempo
($\Delta t=\SI{e-5}{s}$): el frente de onda ni enciende (con el estímulo nominal)
ni conduce (ni siquiera con un estímulo más fuerte), mientras que \gpode\ enciende
y propaga en la misma malla y el mismo $\Delta t$. Esta nota muestra, con tres
corridas controladas de discriminación y un argumento matemático breve, que esto
\emph{no} es un error de código, ni un artefacto del esquema temporal
(Crank--Nicolson), ni un error de escala en $\Delta t$. Es una limitación
intrínseca de precisión al evaluar una cinética iónica fuertemente no lineal en el
potencial de membrana \emph{promediado en la celda}: por la desigualdad de Jensen,
la activación de sodio evaluada en el promedio de celda sub-estima el promedio
puntual, debilitando la corriente regenerativa $I_{\mathrm{Na}}$ en el
\emph{upstroke} (fase de despolarización rápida) abrupto. Subir el orden de
reconstrucción (\code{p2}, \code{p3}) no ayuda, porque el orden solo afecta cómo
se interpolan los \emph{states} ya calculados, no el potencial que manejó el ODE.
\end{abstract}

% ================================================================== %
\section{Ecuaciones de gobierno y los dos modos de integración de \emph{states}}
% ================================================================== %

El problema monodominio resuelto para el potencial transmembrana $\Vm$ es
\begin{equation}
  \chi C_m \,\frac{\partial \Vm}{\partial t}
  \;=\; \divg\!\left(\Dten\,\grad \Vm\right)
        \;-\; \Iion(\Vm,\states) \;+\; \Iext ,
  \qquad
  \frac{\partial \states}{\partial t} \;=\; \mathbf{g}(\Vm,\states),
  \label{eq:monodomain}
\end{equation}
donde $\states$ agrupa las $N$ variables de compuerta (\emph{gating}) y
concentraciones iónicas del modelo celular ten Tusscher--Noble--Noble--Panfilov
(TNNP), $\mathbf{g}$ es su lado derecho rígido, y $\Iion=\phi(\Vm,\states)$ es la
corriente iónica total (algebraica). El término fuente de reacción de
\eqref{eq:monodomain} se integra sobre cada celda $\Omegac$ con una cuadratura de
alto orden en los puntos $\{x_q\}$ y pesos normalizados $\{w_q\}$
($\sum_q w_q=1$; los pesos LRE ya llevan las fracciones de volumen sub-celda, así
que la cuadratura devuelve directamente el \emph{promedio} de celda),
\begin{equation}
  \overline{\Iion}_{,\,c}
  \;=\; \frac{1}{|\Omegac|}\int_{\Omegac} \Iion \,\mathrm{d}\Omega
  \;\approx\; \sum_{q} w_q\,\phi\!\big(\Vm(x_q),\,\states(x_q)\big).
  \label{eq:quad}
\end{equation}
El potencial de membrana en los puntos de Gauss, $\Vm(x_q)$, se obtiene en ambos
modos por la misma reconstrucción de alto orden del campo centrado en celda. Los
dos modos difieren \emph{solo} en cómo se obtiene el vector de \emph{states}
$\states(x_q)$ que entra en \eqref{eq:quad}.

\paragraph{\gpode\ (heredado).}
El ODE rígido se integra \emph{independientemente en cada punto de Gauss},
manejado por el potencial reconstruido localmente:
\begin{equation}
  \states_q^{\,n+1}
  \;=\; \mathrm{ODEsolve}\!\big(\states_q^{\,n},\,\Vm(x_q);\,\Delta t\big),
  \qquad
  \states(x_q) \equiv \states_q^{\,n+1}.
  \label{eq:gp}
\end{equation}
Costo: $N_q$ integraciones rígidas por celda.

\paragraph{\ccrec\ (nuevo).}
El ODE se integra \emph{una vez por celda}, manejado por el potencial de centro de
celda $\Vbarc \equiv \Vm(x_c)$, y los \emph{states} centrados en celda resultantes
se reconstruyen luego en los puntos de Gauss por una LRE de orden $p$:
\begin{equation}
  \states_c^{\,n+1}
  \;=\; \mathrm{ODEsolve}\!\big(\states_c^{\,n},\,\Vbarc;\,\Delta t\big),
  \qquad
  \states(x_q) \equiv \widehat{\states}^{\,(p)}(x_q)
        = \mathcal{R}^{(p)}\!\big[\{\states_{c'}^{\,n+1}\}_{c'\in\text{stencil}}\big](x_q).
  \label{eq:cc}
\end{equation}
Costo: una integración rígida por celda (la motivación: las integraciones TNNP
limitadas por $I_{\mathrm{Na}}$ dominan el tiempo de corrida, así que $1$ vs $N_q$
ODEs por celda es un ahorro grande).

El contraste esencial ya es visible en \eqref{eq:gp}--\eqref{eq:cc}: en \gpode\ el
mapa no lineal $\Vm\!\mapsto\!\states$ se aplica al $\Vm(x_q)$ \emph{puntual}; en
\ccrec\ se aplica al \emph{promedio de celda} $\Vbarc$, y solo el resultado (ya
calculado) se interpola.

% ================================================================== %
\section{Síntoma observado}
% ================================================================== %

Benchmark: losa 2-D, hexaedros estructurados, $\Delta x=\SI{0.5}{mm}$
($40\times16$), TNNP, una caja de estímulo en la esquina de \SI{1.5}{mm} con
intensidad \SI{50000}{A/m^3} durante \SI{2}{ms}, umbral de activación \SI{0}{mV},
Crank--Nicolson, acoplamiento Picard, RKF45 para los \emph{states}. Con \ccrec\ a
orden \code{p1} (``CCp1''):
\begin{itemize}[nosep]
  \item a $\Delta t=\SI{e-4}{s}$ la activación enciende y propaga
        (los puntos de la diagonal se activan en secuencia);
  \item a $\Delta t=\SI{e-5}{s}$ \emph{nada} se activa --- toda muestra de la
        diagonal, incluido el propio sitio del estímulo, permanece en
        $t_{\mathrm{act}}=0$.
\end{itemize}
Que un paso de tiempo más fino (más preciso) produzca \emph{ninguna} activación es
la firma contra-intuitiva que hay que explicar.

% ================================================================== %
\section{Experimentos de discriminación}
% ================================================================== %

Tres corridas controladas aíslan la causa. Todas usan la misma malla
($40\times16$, $\Delta x=\SI{0.5}{mm}$), $\Delta t=\SI{e-5}{s}$, tiempo final
\SI{10}{ms} (suficiente: el sitio del estímulo dispara a ${\sim}\SI{1.3}{ms}$
cuando se gatilla). ``Ignición'' significa que el sitio del estímulo (índice $0$
de la diagonal) cruza el umbral; ``conducción'' significa que el frente avanza a
lo largo de la diagonal.

\begin{table}[h]
\centering
\footnotesize
\begin{tabular}{@{}l l c c p{5.2cm}@{}}
\toprule
Corrida & Configuración & Ignición & Conducción & Evidencia \\
\midrule
A & \gpode\ (\code{states=NO}), CN            & \textbf{sí} & \textbf{sí}
  & pto.\ $0$ @ \SI{1.20}{ms}; ptos.\ $0$--$4$ activos, frente ${\sim}\SI{4.5}{mm}$ \\
--- & \ccrec\ CCp1, CN (original)             & no  & ---
  & todos $t_{\mathrm{act}}=0$; $\max \Vm=\SI{-27}{mV}$ \\
B & \ccrec\ CCp1, \textbf{backwardEuler}      & no  & ---
  & todos $t_{\mathrm{act}}=0$; $\max \Vm=\SI{-27}{mV}$ \\
C & \ccrec\ CCp1, CN, estímulo $\times 4$     & \textbf{sí} & no
  & pto.\ $0$ @ \SI{0.60}{ms}; solo ptos.\ $0$--$2$, el frente se estanca ${\sim}\SI{2.3}{mm}$ \\
\bottomrule
\end{tabular}
\caption{Corridas de discriminación a $\Delta t=\SI{e-5}{s}$. \gpode\ enciende y
conduce; \ccrec\ no hace ninguna de las dos con el estímulo nominal, e incluso un
estímulo $\times 4$ solo enciende sin sostener la conducción.}
\label{tab:disc}
\end{table}

\noindent Las tres corridas descartan los sospechosos ``obvios'' y localizan la causa:
\begin{description}[style=nextline,leftmargin=1.2em,nosep]
  \item[No es el esquema temporal.] La Corrida~B reemplaza Crank--Nicolson por el
    backward Euler totalmente disipativo y aun así falla de forma idéntica
    ($\max\Vm=\SI{-27}{mV}$). Crank--Nicolson queda exonerado.
  \item[No es un error de escala en $\Delta t$.] El esquema $\theta$ es
    $\Delta t$-consistente
    ($\mathsf{A}=\mathsf{M}/\Delta t-\theta\mathsf{L}$,
    $\mathsf{B}=\mathsf{M}/\Delta t+(1-\theta)\mathsf{L}$, la fuente ponderada por
    $\theta,(1-\theta)$ sin factor $\Delta t$ espurio), y el estímulo se aplica
    sobre una ventana física fija $[t_0,t_0+\SI{2}{ms}]$ independiente de
    $\Delta t$. Ambos $\Delta t$ inyectan la misma carga.
  \item[\emph{Sí} es el modo de \emph{states}.] La Corrida~A es idéntica al caso
    que falla salvo que los \emph{states} se resuelven en los puntos de Gauss
    (\gpode). Enciende a \SI{1.20}{ms} y conduce. El único ingrediente cambiado es
    \eqref{eq:gp} vs \eqref{eq:cc}.
\end{description}

% ================================================================== %
\section{Causa raíz: promediar una cinética no lineal}
% ================================================================== %

\subsection{La activación de sodio es convexa en el pie del \emph{upstroke}}

La despolarización rápida la lleva la corriente de sodio
$I_{\mathrm{Na}}=g_{\mathrm{Na}}\,m^{3}h\,j\,(\Vm-E_{\mathrm{Na}})$, cuya
compuerta de activación relaja hacia el sigmoide empinado
\begin{equation}
  m_\infty(\Vm)=\Big[1+\exp\!\big((V_{1/2}-\Vm)/k\big)\Big]^{-1},
  \qquad V_{1/2}\approx\SI{-40}{mV},\; k\approx\SI{6}{mV}.
\end{equation}
Por debajo de su inflexión ($\Vm<V_{1/2}$) --- exactamente el ``pie''
sub-umbral donde se decide la ignición --- $m_\infty$ es \emph{convexa}. Sobre una
celda $\Omegac$ que abarca el frente de onda, $\Vm(x)$ va desde el reposo en el
lado de atrás hasta fuertemente despolarizado en el lado de adelante.

\subsection{Desigualdad de Jensen: \ccrec\ sub-dispara la corriente reactiva}

Sea $\avg{\cdot}=\sum_q w_q(\cdot)$ (con $\sum_q w_q=1$) el promedio de
cuadratura de celda, de modo que $\Vbarc=\avg{\Vm}$. \gpode\ evalúa la cinética
puntualmente y promedia el \emph{resultado}; \ccrec\ promedia la \emph{entrada}
primero. Para la $m_\infty$ convexa,
\begin{equation}
  \underbrace{\avg{\,m_\infty(\Vm)\,}}_{\text{\gpode\ (efectivo)}}
  \;\;\ge\;\;
  \underbrace{m_\infty\!\big(\avg{\Vm}\big)}_{\text{\ccrec}}
  \;=\; m_\infty(\Vbarc).
  \label{eq:jensen}
\end{equation}
Como $I_{\mathrm{Na}}\propto m^{3}$, la brecha en \eqref{eq:jensen} se amplifica al
cubo. \ccrec\ por lo tanto \emph{sub-estima} sistemáticamente la fracción de sodio
abierto, y por ende la corriente entrante (despolarizante), justo donde $\Vm$ varía
más dentro de una celda --- en el borde del estímulo y a lo largo del
\emph{upstroke}. En \gpode\ los puntos de Gauss situados en el lado despolarizado
de la celda corren cada uno su propio ODE, ven un $\Vm(x_q)$ local alto, abren sus
canales de sodio, y aportan una corriente entrante fuerte a \eqref{eq:quad}; esos
``hot spots'' (puntos calientes sub-celda) son exactamente lo que \ccrec\ promedia
y elimina antes de que la cinética siquiera los vea.

\subsection{Dos modos de falla}

\begin{enumerate}[leftmargin=1.4em,nosep]
  \item \textbf{Falla de ignición.} Con el estímulo nominal, el sitio se despolariza
    solo hasta $\Vm\approx\SI{-27}{mV}$ (Corridas~orig/B) y relaja de vuelta: la
    $I_{\mathrm{Na}}$ debilitada nunca se vuelve regenerativa, así que el
    \emph{upstroke} de todo-o-nada no dispara. Un estímulo $\times 4$ (Corrida~C)
    fuerza a $\Vm$ a cruzar el umbral en el sitio ($\SI{0.60}{ms}$), confirmando
    que el promediado dejó al estímulo nominal \emph{efectivamente sub-umbral}.
  \item \textbf{Bloqueo de conducción.} Aun ya encendido (Corrida~C), el frente se
    estanca luego de $2$--$3$ celdas. La propagación requiere que cada celda que se
    activa aporte suficiente carga entrante (\emph{source}, fuente) para llevar a
    las celdas aguas abajo (\emph{sink}, sumidero) al umbral --- el ``factor de
    seguridad'' cardíaco. Por \eqref{eq:jensen}, \ccrec\ baja la fuente justo en el
    frente donde se necesita, llevando el factor de seguridad por debajo de uno: un
    desajuste \emph{source--sink} (fuente--sumidero) inducido numéricamente, es
    decir bloqueo de conducción.
\end{enumerate}

% ================================================================== %
\section{Por qué el paso de tiempo grueso ``funcionaba''}
% ================================================================== %

A $\Delta t=\SI{e-4}{s}$ el paso implícito avanza la despolarización en un solo
salto grande: la fuente de reacción/estímulo actúa sobre un incremento grande
$\Delta t/\mathsf{M}$, produciendo un \emph{over-shoot} (sobre-disparo) transitorio
de $\Vbarc$ que puede exceder el umbral aun con la $I_{\mathrm{Na}}$ sub-estimada.
A medida que $\Delta t\to0$ la trayectoria calculada converge a la \emph{verdadera}
(precisa), que bajo la corriente debilitada de \ccrec\ permanece sub-umbral --- así
que la deficiencia, enmascarada por el sobre-disparo de paso grueso, queda
expuesta. En consecuencia, la activación de CCp1 a $\Delta t=\SI{e-4}{s}$ es un
artefacto de sub-resolución temporal y sus tiempos de activación no deben confiarse;
conviene verificar que coinciden con \gpode\ al mismo $\Delta t$.

% ================================================================== %
\section{Por qué un orden de reconstrucción más alto no ayuda}
% ================================================================== %

Los modos \code{p2} y \code{p3} cambian solo el operador $\mathcal{R}^{(p)}$ en
\eqref{eq:cc} --- cómo se interpolan a los puntos de Gauss los \emph{states}
centrados en celda $\states_c^{\,n+1}$. \emph{No} cambian el hecho de que
$\states_c^{\,n+1}$ se produjo integrando el ODE con la entrada de promedio de
celda $\Vbarc$. El sesgo en \eqref{eq:jensen} se comete \emph{dentro} de la
resolución del ODE, antes de cualquier reconstrucción; una interpolación de orden
más alto solo refina una cantidad ya sesgada. Empíricamente y por construcción,
CCp2/CCp3 heredan la misma falla de ignición y de conducción que CCp1 en este
frente.

% ================================================================== %
\section{Cuándo \ccrec\ es válido, y recomendación}
% ================================================================== %

\eqref{eq:jensen} es una igualdad cuando $\Vm$ es (casi) lineal a lo largo de la
celda, es decir cuando la variación sub-celda es pequeña comparada con la escala de
curvatura de $m_\infty$. Esta es la condición
\begin{equation}
  h_{\text{cell}} \;\ll\; w_{\text{front}},
\end{equation}
con el tamaño de celda mucho menor que el ancho del \emph{upstroke}. El
\emph{upstroke} cardíaco tiene ${\sim}\SI{0.5}{}$--$\SI{1}{mm}$ de ancho; a
$h=\SI{0.5}{mm}$ una celda abarca todo el frente y \ccrec\ falla, mientras que
$h\lesssim\SI{0.1}{mm}$ lo resolvería --- a un conteo de celdas que erosiona la
propia ventaja de costo que \ccrec\ debía proveer.

\begin{itemize}[leftmargin=1.4em,nosep]
  \item \textbf{Propagación con un modelo celular rígido (Niederer/TNNP):} usar
    \gpode\ (\code{stateIntegrationMode gaussPointODE}, o \code{states="NO"} en el
    barrido de tríadas). Resuelve la cinética sub-celda y enciende y conduce.
  \item \textbf{Campos de reacción suaves} (p.\,ej.\ el caso manufacturado MMS del
    solver FDA, donde $\Vm$ varía poco dentro de una celda): \ccrec\ es preciso y
    entrega el ahorro de ODE $N_q\!\to\!1$ que se propone.
  \item \textbf{Híbrido posible (trabajo futuro):} integrar el ODE en los puntos de
    Gauss solo donde $\lvert\grad\Vm\rvert$ es grande (el frente móvil) y en el
    centro de celda con reconstrucción en el resto (reposo/meseta). Esto
    recuperaría la mayor parte del ahorro preservando la precisión en el frente.
\end{itemize}

\paragraph{Conclusión.} \ccrec\ no está roto y su implementación es correcta;
simplemente cambia resolución sub-celda de la cinética iónica por velocidad. Ese
canje es seguro para reacciones suaves pero inválido para el \emph{upstroke}
empinado y auto-sostenido de un potencial de acción cardíaco real, donde promediar
el potencial de membrana antes de la cinética de sodio suprime la ignición y
bloquea la conducción. Para el benchmark de Niederer, \gpode\ es el modo apropiado.

\end{document}
